\subsubsection{Linear Differential Operator}

\paragraph{The Formal Linear Differential Operator}

We define the differentian opeartor as:

\begin{definition}[Differentiation Opeartor]Let $D:L^2[a,b]\to L^2[a,b]$ be the mapping defined by:
\begin{equation}
D(u)=\frac{d}{dx}u
\end{equation}
for all $u\in L^2[a,b]$

\end{definition}Now we can define the \textit{formal linear operator} by some opeartor algebra:

\begin{definition}[Formal linear differential operator]Let $p(x)$ be any polynomial over $\mathbb{C}$ of positive degree:
\begin{equation}
p(x)=a_\nu x^\nu
\end{equation}
where $a_{\nu}$ are functions of $x$. Define $L:=p(D)$ as the \textbf{formal linear differential operator} that correspond to the \textbf{auxillary polynomial $p$}:
\begin{equation}
L=a_\nu D^{\nu}
\end{equation}
The \textbf{order} of $L$ is given by the order of $p$.

\end{definition}We call it \textit{formal} since we want to contrast with \textit{concrete} linear differential opeartor which is a formal opeartor endorsed with boundary conditions(This convention is used in \textit{Mathematical Physics by Michael Stone and Paul Goldbart}). We can think of the difference between them as:

\begin{itemize}
\item The \textbf{nullspace of formal opeartor} is the space of general solution of homogenous diffferential equation
\item The \textbf{nullspace of concrete operator} is the space of particular solution to homogenous differential equation with boundary conditions
\end{itemize}

We'll explore the nullspace of formal operators with a bit more detail now.

\subparagraph{Nullspace of formal operator}

We first observe something very nice about solutions to homogenous differential equations: they belong to $C^{\infty}$(infinitely differentiable).

\subparagraph{The Bootstrapping Lemma}

\begin{lemma}Let $L$ be a formal operator of order $n$. Suppose the coefficients of the auxillary polynomial are $k$ times continously differentiable($a_{\nu}\in C^k$) then $\forall y\in \mathrm{Null}(L)$, $y$ is at least $k+n$ times contiously differentiable.

\end{lemma}\begin{proof}Suppose $y\in \mathrm{Null}(L)$, note trivially $y\in C^n$. Then:
\begin{equation}
a_0y+a_1y^{(1)}+a_2y^{(2)}...+a_n y^{(n)} &=0 \\
y^{(n)} &= -(b_0y+b_1y^{(1)}+b_2y^{(2)}...\underbrace{b_{n-1}{y}^{n-1}}_{\in C^{\min(k, 1)}})
\end{equation}
Where $b_{\nu}=a_{\nu}/a_{n}$. Now $b_{\nu}\in C^k$ by closedness under polynomial division. Also $y^{(\nu)} \in C^{n -\nu}$ hence the R.H.S is at least one time continously differentiable(the bottle neck being $a_{n -1}y^{(n -1)} \in C^{\min(k, 1)}$). Therefore $y^{(n)}\in C^1$. But this implies $y\in C^{n+1}$ hence we repeat the above argument to get $y^{(\nu)}\in C^{n+1 -\nu}$. The R.H.S has the bottleneck $a_{n -1}y^{(n -1)}\in C^{\min(k, n+1 -(n -1))}=C^{\min(k, 2)}$. We can repeat this procedure $m$ times and yield that $y^{(n)}\in C^{\min(k, m)}$. This means the best we can get is $y^{(n)}\in C^k$ or $y\in C^{n+k}$(after which the coeficient becomes the bottleneck).

\end{proof}\begin{corollary}If the coefficients are in $C^{\infty}$ then $y\in \mathrm{Null}(L)$ are in $C^{\infty}$.

\end{corollary}\subparagraph{Nullspace of $L$ is an $n$-dimensional subspace of $C^{\infty}$}

\paragraph{Linear Independence: Wronskian}

\paragraph{Solving $Ly=0$ with constant coefficients}

\paragraph{$Ly=0$ with general $a_\nu(x)$}